\rSec3[box.class]{Class template box}

\begin{codeblock}
namespace std {
  template<class T>
  class @\libglobal{box}@ {
  public:
    // \ref{box.ops}, operations
    template<class F> constexpr auto apply(F&& f) const;
    template<class F> constexpr auto chain(F&& f) const;
    template<class F> constexpr auto local(F&& f) const;
    constexpr auto member() const;
    template<class F> constexpr auto prose(F&& f) const;
    template<class F> constexpr auto triple(F&& f) const;
    template<class F> constexpr auto joined(F&& f) const;

  private:
    T* @\exposid{val}@ = nullptr; // \expos
  };
}
\end{codeblock}

\pnum
The leading aliases that a static_assert names are introduced by one derived "Let" paragraph ahead of the Mandates, and leave the Effects code (issue \#404, decision let-binding-wording).

\rSec3[box.ops]{Operations}

\indexlibrarymember{apply}{box}%
\begin{itemdecl}
template<class F>
  constexpr auto apply(F&& f) const;
\end{itemdecl}

\begin{itemdescr}
\pnum
Let \tcode{U} be \tcode{invoke_result_t<F, T\&>}.

\pnum
\mandates
\tcode{is_object_v<U>} is \tcode{true}.

\pnum
\effects
Equivalent to:
\begin{codeblock}
return invoke(std::forward<F>(f), *@\exposid{val}@);
\end{codeblock}
\end{itemdescr}

\indexlibrarymember{chain}{box}%
\begin{itemdecl}
template<class F>
  constexpr auto chain(F&& f) const;
\end{itemdecl}

\begin{itemdescr}
\pnum
Let \tcode{U} be \tcode{invoke_result_t<F, T\&>} and \tcode{V} be \tcode{remove_cvref_t<U>}.

\pnum
An alias the asserted one names is introduced too, in source order.

\pnum
\mandates
\tcode{is_object_v<V>} is \tcode{true}.

\pnum
\effects
Equivalent to:
\begin{codeblock}
return invoke(std::forward<F>(f), *@\exposid{val}@);
\end{codeblock}
\end{itemdescr}

\indexlibrarymember{local}{box}%
\begin{itemdecl}
template<class F>
  constexpr auto local(F&& f) const;
\end{itemdecl}

\begin{itemdescr}
\pnum
An alias only the body uses stays in the code.

\pnum
\mandates
\tcode{is_object_v<T>} is \tcode{true}.

\pnum
\effects
Equivalent to:
\begin{codeblock}
using W = invoke_result_t<F, T&>;

return W(invoke(std::forward<F>(f), *@\exposid{val}@));
\end{codeblock}
\end{itemdescr}

\indexlibrarymember{member}{box}%
\begin{itemdecl}
constexpr auto member() const;
\end{itemdecl}

\begin{itemdescr}
\pnum
Let \tcode{U} be \tcode{decltype((\exposid{val}))}.

\pnum
\mandates
\tcode{is_const_v<U>} is \tcode{false}.

\pnum
\effects
Equivalent to:
\begin{codeblock}
return *@\exposid{val}@;
\end{codeblock}
\end{itemdescr}

\indexlibrarymember{prose}{box}%
\begin{itemdecl}
template<class F>
  constexpr auto prose(F&& f) const;
\end{itemdecl}

\begin{itemdescr}
\pnum
Let \tcode{U} be \tcode{invoke_result_t<F, T\&>}.

\pnum
\tcode{U} names the result of calling \tcode{f}.

\pnum
\mandates
\tcode{is_object_v<U>} is \tcode{true}.

\pnum
\effects
Equivalent to:
\begin{codeblock}
return invoke(std::forward<F>(f), *@\exposid{val}@);
\end{codeblock}
\end{itemdescr}

\indexlibrarymember{triple}{box}%
\begin{itemdecl}
template<class F>
  constexpr auto triple(F&& f) const;
\end{itemdecl}

\begin{itemdescr}
\pnum
Let \tcode{U} be \tcode{invoke_result_t<F, T\&>}, \tcode{V} be \tcode{remove_cvref_t<U>}, and \tcode{W} be \tcode{remove_cvref_t<V>}.

\pnum
Three aliases are bound in one sentence, the last after ", and".

\pnum
\mandates
\tcode{is_object_v<W>} is \tcode{true}.

\pnum
\effects
Equivalent to:
\begin{codeblock}
return invoke(std::forward<F>(f), *@\exposid{val}@);
\end{codeblock}
\end{itemdescr}

\indexlibrarymember{joined}{box}%
\begin{itemdecl}
template<class F>
  constexpr auto joined(F&& f) const;
\end{itemdecl}

\begin{itemdescr}
\pnum
Let \tcode{U} be \tcode{invoke_result_t<F, T\&>} and \tcode{V} be \tcode{remove_cvref_t<U>}.

\pnum
\mandates
\tcode{is_object_v<V>} is \tcode{true}.

\pnum
\effects
Equivalent to:
\begin{codeblock}
return invoke(std::forward<F>(f), *@\exposid{val}@);
\end{codeblock}
\end{itemdescr}
